\source{https://aclanthology.org/2026.acl-demo.27/；逐页补充原文页码与图表编号}
\papersection{Hindsight}{保留、召回与反思的结构化记忆}{讲解人待填写}

\begin{frame}{研究问题与任务范围}
    \large\key{待补充：论文解决的具体问题}\par
    \vspace{5mm}\normalsize
    \begin{itemize}\setlength{\itemsep}{4mm}
        \item 任务、输入与预期输出是什么？
        \item 既有方法在哪种情形下失败？
        \item 论文提出什么可检验的改进目标？
    \end{itemize}
    \note{建议约 45 秒。用一个具体问题或案例说明动机，避免用过长的背景综述占用时间。}
\end{frame}

\begin{frame}{核心方法与整体流程}
    \begin{columns}[T,onlytextwidth]
        \column{.54\textwidth}
        \textbf{核心流程图：待替换}\par\vspace{5mm}
        {\color{muted}放入原论文核心图或可编辑重绘图。\par
        保留关键标签，标明输入、操作与输出。}
        % \includegraphics[width=\linewidth,height=.50\textheight,keepaspectratio]{images/paper-method.pdf}
        \column{.42\textwidth}
        \textbf{需要解释的机制}\par\vspace{4mm}
        \begin{enumerate}\setlength{\itemsep}{4mm}
            \item 待补充：信息如何进入系统。
            \item 待补充：核心操作与设计原因。
            \item 待补充：结果如何被后续使用。
        \end{enumerate}
    \end{columns}
    \note{建议约 90 秒。先沿流程讲一遍，再解释核心机制。需要复杂图示时拆页，不把整篇系统图缩成不可读的小图。}
\end{frame}

\begin{frame}{关键机制与方法细节}
    \[
        \text{输入}\quad\longrightarrow\quad
        \text{核心操作}\quad\longrightarrow\quad\text{输出}
    \]
    \vspace{4mm}
    \begin{itemize}\setlength{\itemsep}{4mm}
        \item 待补充：关键公式、结构或操作，说明符号与含义。
        \item 待补充：一个小例子，解释机制怎样改变处理结果。
        \item 待补充：与基线的具体差别，以及依赖的条件。
    \end{itemize}
    \note{建议约 90 秒。聚焦最重要的一项设计；公式要解释符号和直觉。其余实现细节可放到备用页。}
\end{frame}

\begin{frame}{实验设置与评价范围}
    \begin{tabularx}{\linewidth}{@{}p{29mm}Y@{}}
        \toprule
        项目 & 依据原论文补充 \\
        \midrule
        任务与数据 & 待填写：数据来源、划分与任务范围 \\
        模型与基线 & 待填写：模型角色、基线与主要配置 \\
        评价指标 & 待填写：指标含义、评价者与计分方式 \\
        预算与成本 & 待填写：上下文预算、构建或推理开销 \\
        \bottomrule
    \end{tabularx}
    \aside{实验条件解释到足以理解结果；训练、优化与测试时可用信息分别说明。}
    \note{建议约 60 秒。不要逐项念参数，重点解释影响结果理解和可比性的条件。}
\end{frame}

\begin{frame}{关键结果与证据}
    \textbf{待补充：这一页数据支持的一项结论}\par\vspace{5mm}
    \begin{tabularx}{\linewidth}{@{}Yrrr@{}}
        \toprule
        配置 & 主指标 & 关键分项 & 成本 \\
        \midrule
        基线 & 待填 & 待填 & 待填 \\
        本文方法 & 待填 & 待填 & 待填 \\
        关键消融 & 待填 & 待填 & 待填 \\
        \bottomrule
    \end{tabularx}\par
    \vspace{5mm}
    \key{待补充：提升出现在哪些条件下，哪里没有改善}
    \aside{用原表数值或图表替换占位，写清单位；论文结果与本组实验分开。}
    \note{建议约 90 秒。解释一个关键对照或消融，不逐格朗读主表；图表较复杂时可分成结果和消融两页。}
\end{frame}

\begin{frame}{局限、阅读判断与后续问题}
    \begin{columns}[T,onlytextwidth]
        \column{.47\textwidth}
        \textbf{论文支持的结论}\par\vspace{4mm}
        \begin{itemize}\setlength{\itemsep}{3mm}
            \item 待补充：最重要的贡献与证据。
            \item 待补充：尚未解决的局限。
        \end{itemize}
        \column{.47\textwidth}
        \textbf{个人阅读思考}\par\vspace{4mm}
        \begin{itemize}\setlength{\itemsep}{3mm}
            \item 待补充：最值得借鉴的设计。
            \item 待补充：需要进一步验证的问题。
        \end{itemize}
    \end{columns}
    \vspace{5mm}
    \aside{与后续讲解的衔接：待各成员填写完成后统一整理。}
    \note{建议约 75 秒。个人判断与论文结论分别表达，并用一句话交接给下一位成员。连同章节页，本模块约 8 分钟。}
\end{frame}
